\chapter{Спецификация версии III}

Версия III начинается не с новой численной формулы, а с одного вопроса:
существует ли действие, которое до обращения к наблюдаемым одновременно
фиксирует меру, embedding-карты коллективных координат и не менее двух
normalization-sensitive секторов.

\section{Замороженные входы}

На первом этапе сохраняются только математические данные, уже отделённые от
физических предсказаний:
\begin{equation}
 K=\mathbb{RP}^3\times S^1,
 \qquad
 \Vol(\mathbb{RP}^3)=\pi^2,
 \qquad
 \ell(S^1)=2\pi,
 \qquad
 \ell(\gamma)=\pi.
\end{equation}
Допустимы projector ranks, integral cohomology classes и представления,
если они заданы до сравнения с массами. Не допускается использование
\(\alpha^{-1}\), \(m_\tau\), neutrino splittings или EW/QCD residuals для
выбора коэффициентов действия.

\begin{protocol}[Двухсекторный Definition of Done]
Кандидат считается первым положительным parent action только если:
\begin{enumerate}
 \item полный набор полей, симметрий, мер и boundary terms объявлен заранее;
 \item одна и та же Hessian-метрика задаёт neutrino и charged-lepton
 collective maps без независимых sector coefficients;
 \item embedding каждой collective coordinate фиксирован калибровочной
 симметрией, интегральным периодом, канонической скобкой или другим
 field-redefinition-covariant принципом;
 \item vacuum equations выбирают branch и family axis либо доказывают их
 физическую несущественность;
 \item действие проходит третий тест, не использованный при его построении;
 \item все машинные расчёты воспроизводимы отдельным аудитом.
\end{enumerate}
\end{protocol}

\section{Порядок работы}

Исследование разбивается на последовательные kill-gates:
\begin{enumerate}
 \item role-graded Hessian: проверить числовую совместимость одной Hodge
 metric при разных типах коллективных координат;
 \item embedding gate: определить, фиксированы ли масштабы этих координат
 симметрией, а не выбором базиса;
 \item vertex covariance: вывести совместное преобразование stiffness и
 interaction vertices;
 \item vacuum selector: проверить branch, axis и rank-one kernel;
 \item independent sector: вычислить заранее выбранный третий observable.
\end{enumerate}

Каждый отрицательный этап закрывает только объявленный класс. Переход к
следующему классу требует нового структурного принципа, а не дополнительного
числового коэффициента.